\section{Introduction}
\label{sec:introduction}

Concentration gradient generators are indispensable in microfluidics for a wide range of biological and chemical applications, including cell chemotaxis studies, drug screening, and nanoparticle synthesis~\cite{whitesides2006}.
The classic ``Christmas tree'' network introduced by Jeon \emph{et~al.} (2000) demonstrated that multiple fluid streams, repeatedly split and recombined in a planar network of microchannels, can produce a smooth, predictable concentration profile at the outlet~\cite{jeon2000}.
This design has been widely adopted because it operates entirely on laminar flow and diffusion, requiring no moving parts.
However, the Christmas tree relies on a fixed, hand-drawn branching topology: the network geometry is prescribed \emph{a~priori}, and the resulting outlet profile is essentially a binomial distribution -- a single family of monotonic, symmetric concentration profiles.
Modifying the profile shape requires cumbersome, trial-and-error re-dimensioning of each channel segment, and the design offers no guarantee of optimality in terms of pressure drop, mixing length, or channel volume.

\subsection{Topology optimisation of fluidic systems}
A powerful alternative to manual geometric design is topology optimisation, a computational method that automatically distributes material within a design domain to extremise a user-defined objective function.
The seminal work of Borrvall and Petersson (2003) established the density-based approach for Stokes flow by introducing an inverse permeability that penalises solid regions, thereby transforming the shape design problem into a continuous material distribution problem solvable with gradient-based optimisers~\cite{borrvall2003}.
Their formulation has become the foundation for a growing body of research in fluidic topology optimisation~\cite{gersborg2005,olesen2006,evgrafov2005,challis2009}.

In the two decades since, the topology optimisation community has extended the Borrvall--Petersson framework in several directions directly relevant to the present work.
Poulsen, Aage, Gersborg-Hansen, and Sigmund addressed the scalability of Stokes flow topology optimisation to large-scale three-dimensional problems, demonstrating that the density method can handle industrial-sized design domains when combined with efficient iterative solvers and parallel computing~\cite{aage2008,poulsen2003}.
Kreissl, Pingen, Evgrafov, and Maute introduced multi-objective formulations for dynamically tunable fluidic devices, where the channel layout is designed to accommodate mechanical actuation and fluid--structure interaction~\cite{kreissl2010,kreissl2011}.
Their work also pioneered the application of the level-set method to fluid topology optimisation, providing sharp fluid--solid interfaces and enabling unsteady flow design~\cite{kreissl2012}.

\subsection{Coupling flow with scalar transport}
A critical step toward designing functional microfluidic devices -- rather than merely minimising pressure drop -- is the simultaneous optimisation of the flow field and a transported scalar quantity, such as temperature or concentration.
Alexandersen, Andreasen, Aage, and Sigmund developed a density-based topology optimisation framework for coupled convection problems, applying it to forced and natural convection heat sinks and micropumps~\cite{alexandersen2014,alexandersen2020}.
Makhija and Maute extended level-set topology optimisation to scalar transport problems, using the extended finite element method (XFEM) to capture concentration fields on evolving interfaces~\cite{makhija2015}.
Deng (2018) published the first monograph dedicated to topology optimisation theory for laminar flow, comprehensively covering the density method and level-set method for inverse design of microfluidic components including passive micromixers, electroosmotic micromixers, no-moving-part microvalves, and valveless micropumps~\cite{deng2018}.

\subsection{Open-source and multiscale developments}
Na \emph{et~al.} (2022) provided an open-source topology optimisation framework for passive micromixer design, employing continuous adjoint sensitivity and demonstrating that topology optimisation can generate smooth, manufacturable channel geometries competitive with hand-designed serpentine mixers~\cite{na2022}.
Feppon (2024) introduced a multiscale topology optimisation method for fluid microchannels based on asymptotic homogenisation, enabling the design of spatially modulated channel arrays with enhanced heat and mass transfer performance~\cite{feppon2024}.
A comprehensive review by Li \emph{et~al.} (2023) surveyed the technical progress of fluid topology optimisation over the preceding decade, highlighting the growing adoption of the density method for microfluidic mixers, heat exchangers, and flow distributors~\cite{li2023}.

\subsection{Inverse design of concentration gradient generators}
Recently, deep-learning-based inverse design has been applied to concentration gradient generators.
Yang and Wang (2020) trained a deep neural network on a physics-based component model to predict the design parameters required to achieve a desired concentration profile in complex network topologies~\cite{yang2020}.
While this approach demonstrates the potential of data-driven inverse design, it still operates within the space of pre-conceived network architectures and does not discover novel channel layouts.

\subsection{Knowledge gap and contributions}
Despite these advances, three important gaps remain.
First, density-based topology optimisation has \emph{not} been applied to the inverse design of concentration gradient generators: the simultaneous optimisation of flow and convection-diffusion transport to achieve an \emph{arbitrary user-specified concentration profile at the outlet} is, to the best of our knowledge, not previously reported in the literature.
Second, the existing topology-optimised microfluidic components (micromixers, valves, pumps, heat sinks) have been designed for mixing efficiency or pressure-drop minimisation, not for producing a spatially controlled concentration field.
Third, no open-source, fully reproducible Python pipeline exists that combines the forward model, adjoint sensitivity, and optimiser into a single tool accessible to experimental microfluidics researchers without a commercial software license.

In this paper, we address all three gaps by:
\begin{enumerate}
    \item Formulating a density-based topology optimisation problem for coupled Brinkman--convection-diffusion transport, with a multi-objective cost function that simultaneously minimises viscous dissipation and outlet concentration profile mismatch.
    \item Deriving the continuous adjoint equations and implementing them in FEniCSx, enabling efficient gradient computation with respect to all design variables simultaneously, independently of their number.
    \item Publishing the complete pipeline as an open-source Python package (\texttt{micrograd}), including mesh generation, forward and adjoint solvers, optimisation routines (OC and MMA), and post-processing tools.
    \item Demonstrating that the optimised designs outperform the classic Christmas tree generator by \textbf{99.7\%} in hydraulic resistance (from \(1.77\times10^{11}\) to \(4.48\times10^{8}\)~Pa\(\cdot\)s/m\(^3\)) while maintaining equivalent outlet concentration fidelity
    (RMSE~$=0.098$; dimensionless, concentration normalised to $[0,1]$).
    \item Showing that the framework can generate non-obvious channel topologies for arbitrary target profiles -- linear, sigmoid, double-peak, and stair-step -- that are impossible to produce with fixed tree-shaped generators.
\end{enumerate}

The remainder of this paper is organised as follows.
Section~\ref{sec:model} presents the mathematical model and the adjoint sensitivity analysis.
Section~\ref{sec:methods} describes the numerical implementation and the open-source \texttt{micrograd} package.
Section~\ref{sec:results} presents the results, including mesh convergence, validation against Navier--Stokes, comparison with the Christmas tree, and the gallery of target profiles.
Section~\ref{sec:discussion} discusses the limitations and future research directions, and Section~\ref{sec:conclusion} concludes the paper.

\paragraph{Recent developments (2021--2026).}
Several advances since 2021 widen the design space for computational
inverse design of microfluidic devices and motivate the present work.

\emph{Differentiable CFD.}
The emergence of fully differentiable fluid solvers---frameworks in which
every arithmetic operation is tracked by automatic differentiation so that
exact gradients propagate through the PDE solve---has made continuous
adjoint methods more accessible~\cite{bezgin2023jaxfluids}.
These tools support the kind of gradient-based optimisation loop we employ,
and confirm that adjoint sensitivity remains the method of choice for
PDE-constrained design in the low-Reynolds-number regime.

\emph{Physics-informed and neural-operator surrogates.}
Physics-informed neural networks with hard constraints (hPINNs) have been
applied to topology optimisation as a mesh-free alternative to FEM-based
adjoint methods~\cite{lu2021hpinn,karniadakis2021pinns}.
Neural operators (Fourier-DeepONet and related architectures) have been
used to replace the PDE forward solve inside the optimisation loop,
achieving order-of-magnitude speedups for channel topology
problems~\cite{kou2025neural}.
These approaches are complementary to, rather than competitive with, the
density-based adjoint method presented here: they require large offline
training datasets that are unavailable for the present gradient-generation
problem, where the target profile itself is the design input.

\emph{Topology optimisation for transport.}
Papadopoulos \& S{\"u}li~\cite{papadopoulos2022stokes} provided the first
rigorous finite-element convergence analysis for the Borrvall--Petersson
Stokes topology optimisation problem, establishing strong convergence of
discrete adjoint solutions and underpinning the numerical reliability of
the formulation adopted here.
Liu et al.~\cite{liu2025phasefield} extended density-based topology
optimisation to microfluidic mixers using a phase-field regularisation,
demonstrating 2D and 3D designs for concentration uniformisation.
Critically, none of these works target an \emph{arbitrary prescribed
outlet profile}: they optimise mixing uniformity or flow dissipation, not
the shape of the outlet concentration distribution.
The present study fills this gap.

